% GENERATED FILE. Edit canonical lessons, then run npm run generate:books.
% canonical-source-hash: fnv1a64:9afff6181ed6e119
% canonical-lessons: PA-C120-rumal, PA-C120-savaitar, PA-C120-kot, PA-C120-anguthi, PA-C120-kangan

\chapter{Handkerchief, Sweater, Coat, Ring, Bangle}
\label{ch:pa-handkerchief-sweater-coat-ring-bangle}
\hlchaptermodality{120}

\begin{chapteropening}
\textbf{By the end of this chapter:} \emph{I can say a handkerchief, a sweater, a coat, a ring and a bangle.}

Complete the last lesson of chapter 120: I can say a handkerchief, a sweater, a coat, a ring and a bangle.
\end{chapteropening}

\section[\texorpdfstring{\pa{ਰੁਮਾਲ} (rumāl)}{ਰੁਮਾਲ (rumāl)}]{\pa{ਰੁਮਾਲ} (rumāl) — a handkerchief}

\label{lesson:PA-C120-rumal}



\begin{quote}
Before the new one: say the Punjabi for a scarf, then the Punjabi for a belt.
\end{quote}

\subsection*{You'll want to know: \pa{ਰੁਮਾਲ}}
\textbf{\pa{ਰੁਮਾਲ}} — \emph{rumāl} — \textquotedblleft{}a handkerchief\textquotedblright{}.

\begin{cousinweb}
\textbf{\pa{ਰੁਮਾਲ}} (\emph{rumāl}) is a loan, usually traced to Persian \emph{rūmāl}, \textquotedblleft{}handkerchief\textquotedblright{}. It is built from \emph{rū}, \textquotedblleft{}face\textquotedblright{}, and \emph{māl}, \textquotedblleft{}wiping\textquotedblright{}.
\end{cousinweb}

\begin{grammarlens}[title={What you've built}]
One more word for this chapter.
\end{grammarlens}

\subsection*{Your turn}
\begin{itemize}
\raggedright
  \item \emph{Say it:} \emph{rumāl}
  \item \emph{Say it:} \emph{rumāl}, once more
  \item \emph{Say it:} say \emph{rumāl}
  \item \emph{Recall:} say the Punjabi for a scarf, then the Punjabi for a belt, then say \emph{rumāl} again
\end{itemize}

\begin{tcolorbox}[breakable,colback=teal!4,colframe=teal!35!black,title={Before you move on}]
What does \pa{ਰੁਮਾਲ} mean? (\textquotedblleft{}A handkerchief\textquotedblright{}.) Say it once more.
\end{tcolorbox}
\section[\texorpdfstring{\pa{ਸਵੈਟਰ} (savaiṭar)}{ਸਵੈਟਰ (savaiṭar)}]{\pa{ਸਵੈਟਰ} (savaiṭar) — a sweater}

\label{lesson:PA-C120-savaitar}



\begin{quote}
Before the new one: say the Punjabi for a belt, then the Punjabi for a handkerchief.
\end{quote}

\subsection*{You'll want to know: \pa{ਸਵੈਟਰ}}
\textbf{\pa{ਸਵੈਟਰ}} — \emph{savaiṭar} — \textquotedblleft{}a sweater\textquotedblright{}.

\begin{cousinweb}
\textbf{\pa{ਸਵੈਟਰ}} (\emph{savaiṭar}) is a loan from English \emph{sweater}.
\end{cousinweb}

\begin{grammarlens}[title={What you've built}]
One more word for this chapter.
\end{grammarlens}

\subsection*{Your turn}
\begin{itemize}
\raggedright
  \item \emph{Say it:} \emph{savaiṭar}
  \item \emph{Say it:} \emph{savaiṭar}, once more
  \item \emph{Say it:} say \emph{savaiṭar}
  \item \emph{Recall:} say the Punjabi for a belt, then the Punjabi for a handkerchief, then say \emph{savaiṭar} again
\end{itemize}

\begin{tcolorbox}[breakable,colback=teal!4,colframe=teal!35!black,title={Before you move on}]
What does \pa{ਸਵੈਟਰ} mean? (\textquotedblleft{}A sweater\textquotedblright{}.) Say it once more.
\end{tcolorbox}
\section[\texorpdfstring{\pa{ਕੋਟ} (koṭ)}{ਕੋਟ (koṭ)}]{\pa{ਕੋਟ} (koṭ) — a coat}

\label{lesson:PA-C120-kot}



\begin{quote}
Before the new one: say the Punjabi for a handkerchief, then the Punjabi for a sweater.
\end{quote}

\subsection*{You'll want to know: \pa{ਕੋਟ}}
\textbf{\pa{ਕੋਟ}} — \emph{koṭ} — \textquotedblleft{}a coat\textquotedblright{}.

\begin{cousinweb}
\textbf{\pa{ਕੋਟ}} (\emph{koṭ}) is a loan from English \emph{coat}. Punjabi has an older, unrelated \emph{koṭ}, \textquotedblleft{}fort\textquotedblright{}.
\end{cousinweb}

\begin{grammarlens}[title={What you've built}]
One more word for this chapter.
\end{grammarlens}

\subsection*{Your turn}
\begin{itemize}
\raggedright
  \item \emph{Say it:} \emph{koṭ}
  \item \emph{Say it:} \emph{koṭ}, once more
  \item \emph{Say it:} say \emph{koṭ}
  \item \emph{Recall:} say the Punjabi for a handkerchief, then the Punjabi for a sweater, then say \emph{koṭ} again
\end{itemize}

\begin{tcolorbox}[breakable,colback=teal!4,colframe=teal!35!black,title={Before you move on}]
What does \pa{ਕੋਟ} mean? (\textquotedblleft{}A coat\textquotedblright{}.) Say it once more.
\end{tcolorbox}
\section[\texorpdfstring{\pa{ਅੰਗੂਠੀ} (angūṭhī)}{ਅੰਗੂਠੀ (angūṭhī)}]{\pa{ਅੰਗੂਠੀ} (angūṭhī) — a ring}

\label{lesson:PA-C120-anguthi}



\begin{quote}
Before the new one: say the Punjabi for a sweater, then the Punjabi for a coat.
\end{quote}

\subsection*{You'll want to know: \pa{ਅੰਗੂਠੀ}}
\textbf{\pa{ਅੰਗੂਠੀ}} — \emph{angūṭhī} — \textquotedblleft{}a ring\textquotedblright{}.

\begin{cousinweb}
\textbf{\pa{ਅੰਗੂਠੀ}} (\emph{angūṭhī}) is a tadbhava, a word inherited and reshaped by sound change. It comes from \emph{*aṅguṣṭhiya}, \textquotedblleft{}of the thumb\textquotedblright{}, a form CDIAL reconstructs from its descendants (CDIAL 138). It is built on Sanskrit \emph{aṅguṣṭha}, \textquotedblleft{}thumb\textquotedblright{}.
\end{cousinweb}

\begin{grammarlens}[title={What you've built}]
One more word for this chapter.
\end{grammarlens}

\subsection*{Your turn}
\begin{itemize}
\raggedright
  \item \emph{Say it:} \emph{angūṭhī}
  \item \emph{Say it:} \emph{angūṭhī}, once more
  \item \emph{Say it:} say \emph{angūṭhī}
  \item \emph{Recall:} say the Punjabi for a sweater, then the Punjabi for a coat, then say \emph{angūṭhī} again
\end{itemize}

\begin{tcolorbox}[breakable,colback=teal!4,colframe=teal!35!black,title={Before you move on}]
What does \pa{ਅੰਗੂਠੀ} mean? (\textquotedblleft{}A ring\textquotedblright{}.) Say it once more.
\end{tcolorbox}
\section[\texorpdfstring{\pa{ਕੰਗਣ} (kangaṇ)}{ਕੰਗਣ (kangaṇ)}]{\pa{ਕੰਗਣ} (kangaṇ) — a bangle}

\label{lesson:PA-C120-kangan}



\begin{quote}
Before the new one: say the Punjabi for a coat, then the Punjabi for a ring.
\end{quote}

\subsection*{You'll want to know: \pa{ਕੰਗਣ}}
\textbf{\pa{ਕੰਗਣ}} — \emph{kangaṇ} — \textquotedblleft{}a bangle\textquotedblright{}.

\begin{cousinweb}
\textbf{\pa{ਕੰਗਣ}} (\emph{kangaṇ}) is a tadbhava, a word inherited from Sanskrit and reshaped by sound change. It comes from \emph{kaṅkaṇa}, \textquotedblleft{}bracelet\textquotedblright{} (CDIAL 2597).
\end{cousinweb}

\begin{grammarlens}[title={What you've built}]
One more word for this chapter.
\end{grammarlens}

\subsection*{Your turn}
\begin{itemize}
\raggedright
  \item \emph{Say it:} \emph{kangaṇ}
  \item \emph{Say it:} \emph{kangaṇ}, once more
  \item \emph{Say it:} say \emph{kangaṇ}
  \item \emph{Recall:} say the Punjabi for a coat, then the Punjabi for a ring, then say \emph{kangaṇ} again
\end{itemize}

\begin{tcolorbox}[breakable,colback=teal!4,colframe=teal!35!black,title={Before you move on}]
What does \pa{ਕੰਗਣ} mean? (\textquotedblleft{}A bangle\textquotedblright{}.) Say it once more.
\end{tcolorbox}
